% Notas de aula — Capítulo 4: Vapor d'Água
% Curso: Meteorologia Prática (baseado em R. Stull, Practical Meteorology, CC BY-NC-SA 4.0)
% Caixas verdes = conteúdo literal do quadro; linhas verdes = encenação.
\documentclass[11pt]{article}
\usepackage{../comum/preambulo}
\graphicspath{{../figuras/cap04/}}

\title{Capítulo 4 --- Vapor d'Água\\
{\large Notas de aula (roteiro de quadro-negro)}}
\author{Curso de Meteorologia Prática}
\date{}

\begin{document}
\maketitle
\thispagestyle{fancy}

\noindent\textbf{Plano sugerido:} 2 aulas de 2\,h.
\emph{Aula 1:} \S\ref{sec:cc}--\S\ref{sec:variaveis} (Clausius--Clapeyron e o
zoológico de variáveis de umidade).
\emph{Aula 2:} \S\ref{sec:lcl}--\S\ref{sec:balanco} (LCL, bulbo úmido, água
precipitável e balanço de água).

\section{A molécula que comanda o tempo}

\begin{noquadro}[abertura]
\textbf{Cap.~4 --- Vapor d'água}\\[2pt]
``N$_2$, O$_2$, Ar: coadjuvantes. O protagonista é o vapor d'água.''\\[2pt]
por quê:\ (i) carrega $L_v = 2{,}5$\,MJ/kg;\ (ii) maior gás de estufa;\
(iii) é o único que \emph{muda de fase} aqui.
\end{noquadro}

Os três papéis já foram anunciados: o calor latente como bateria térmica
\javimos{3}{calor latente}, o vapor como principal absorvedor de IV
\javimos{2}{efeito estufa}, e a condensação que fará nuvens
\veremos{6}{nuvens} e chuva \veremos{7}{precipitação}. Este capítulo é o
dicionário quantitativo da umidade — e a chave é uma única curva
exponencial, $e_s(T)$.

\section{Pressão de vapor de saturação: Clausius--Clapeyron}\label{sec:cc}

\quadro{Desenhar recipiente fechado com água: setas de evaporação e
condensação; equilíbrio dinâmico define $e_s(T)$. Depois desenhar a curva
$e_s(T)$ — ela fica no quadro a aula inteira.}

Saturação \emph{não} é ``o ar segurar o máximo de água'' (o ar nem
participa!): é o equilíbrio dinâmico entre a taxa de moléculas que escapam
do líquido (cresce explosivamente com $T$) e a taxa que retorna
($\propto e$). A pressão parcial de equilíbrio $e_s$ depende \emph{só} de
$T$ — e de forma violentamente não linear:

\begin{formadif}[a lei que governa a curva]
Da termodinâmica de transições de fase (igualdade dos potenciais de Gibbs
das duas fases na coexistência):
\[ \frac{de_s}{dT} = \frac{L_v\, e_s}{\Re_v\, T^2},\qquad
\Re_v = 461\ \mathrm{J\,kg^{-1}K^{-1}}. \]
Integrando com $L_v$ constante de $(T_0, e_0)$ a $(T, e_s)$:
\[ e_s = e_0 \exp\!\left[\frac{L_v}{\Re_v}\left(\frac{1}{T_0}-\frac{1}{T}\right)\right]. \]
No notebook (Caso~1) integramos a EDO numericamente e comparamos com esta
fórmula e com a do MetPy.
\end{formadif}

\begin{formalivro}[Stull eq.~4.1]
\[ e_s = e_0 \exp\!\left[ 5423\ \mathrm{K}\left(\frac{1}{273{,}15}
 - \frac{1}{T}\right)\right],\qquad e_0 = 0{,}6113\ \mathrm{kPa}, \]
com $L_v/\Re_v = 5423$\,K sobre água líquida (e $L_d/\Re_v = 6139$\,K sobre
gelo — a diferença alimenta o processo de Bergeron
\veremos{7}{processo de Bergeron}).
\end{formalivro}

\begin{noquadro}[os números da curva]
$e_s$ \textbf{dobra a cada $\sim$10\,K}\quad ($\simeq 7\%$ por K)\\[2pt]
$e_s(10\,°$C$) = 1{,}23$\,kPa \quad $e_s(20\,°$C$) = 2{,}34$\,kPa \quad
$e_s(30\,°$C$) = 4{,}24$\,kPa\\[2pt]
$\Rightarrow$ trópicos carregam $3$--$4\times$ mais vapor que latitudes
médias.
\end{noquadro}

\begin{exemplo}[a não linearidade que faz o clima]
$e_s(30\,°\mathrm{C})/e_s(10\,°\mathrm{C}) = 4{,}24/1{,}23 \simeq 3{,}5$.
Ar tropical saturado carrega 3,5$\times$ mais vapor que ar de inverno
ameno — por isso a chuva tropical é mais violenta, e os 7\,\%/K são a taxa
que aparece nas projeções de chuvas extremas com aquecimento global
\veremos{21}{extremos num clima quente}.
\emph{Verificação: $1{,}07^{20} = 3{,}9$ — ordem certa.}
\end{exemplo}

Abaixo de 0\,°C há \emph{duas} curvas: sobre água super-resfriada e sobre
gelo, com $e_s^{gelo} < e_s^{liq}$ (o gelo prende mais as moléculas). Um
ar pode estar subsaturado para a água e supersaturado para o gelo ao mesmo
tempo — guardar essa assimetria: ela fará crescer os cristais às custas
das gotículas \veremos{7}{Bergeron}.

\section{O zoológico de variáveis de umidade}\label{sec:variaveis}

\quadro{Montar a tabela no quadro aos poucos, cada variável com sua
pergunta-guia. Não despejar tudo de uma vez.}

\begin{noquadro}[o zoológico, com as perguntas-guia]
\begin{tabular}{lll}
$e$ & pressão parcial do vapor & quanto vapor há (em kPa)?\\
$r = \dfrac{\varepsilon e}{P-e}$ & razão de mistura & g de vapor por kg de
ar seco?\\
$q \simeq r/(1+r)$ & umidade específica & fração da massa total?\\
$\rho_v = e/\Re_v T$ & umidade absoluta & g de vapor por m$^3$?\\
UR $= 100\,e/e_s$ & umidade relativa & quão perto da saturação?\\
$T_d$: $e_s(T_d)=e$ & ponto de orvalho & até que $T$ esfriar para
saturar?\\
$T_w$ & bulbo úmido & quanto a evaporação consegue esfriar?\\
\end{tabular}
\end{noquadro}

\begin{formalivro}[razão de mistura, Stull eq.~4.4]
\[ r = \frac{\varepsilon\, e}{P - e},\qquad
\varepsilon = \frac{M_{H_2O}}{M_{ar}} = 0{,}622. \]
$r$ é a variável \emph{conservada} em subidas sem condensação — o CPF da
umidade da parcela, par do $\theta$ \javimos{3}{temperatura potencial}.
O $\varepsilon = 0{,}622$ é o mesmo da correção virtual
\javimos{1}{temperatura virtual} (exercício T3 do Cap.~1).
\end{formalivro}

\quadro{Insistir na hierarquia: $e$ e $\rho_v$ mudam com $T$ e $P$; $r$ e
$q$ só mudam se entrar/sair água. UR é a PIOR variável para comparar
massas de ar — depende de $T$! (O ar da Sibéria a $-40\,°$C com UR
$=100\%$ é mais seco que o do Saara.)}

O par conservado $(\theta, r)$ é a impressão digital de uma massa de ar —
é com ele que classificaremos massas polares e tropicais
\veremos{12}{massas de ar} e diagnosticaremos a camada de mistura
\veremos{18}{perfis da camada limite}.

\begin{exemplo}[dia úmido de verão em São Paulo]
$T = 28\,°$C, $T_d = 22\,°$C, $P = 93$\,kPa (altitude da cidade).
$e = e_s(22\,°\mathrm{C}) = 2{,}64$\,kPa;
$r = 0{,}622\times2{,}64/(93-2{,}64) = 18{,}2$\,g/kg;
$\mathrm{UR} = 100\times2{,}64/3{,}78 = 70\,\%$.
\emph{Verificação: 15--20\,g/kg é típico de ar tropical úmido; com
$T_v = T(1+0{,}61r)$ do Cap.~1, esse ar ``vale'' $+3{,}3$\,K de
flutuabilidade.}
\end{exemplo}

\section{Nível de condensação por levantamento (LCL)}\label{sec:lcl}

\quadro{Desenhar parcela subindo num gráfico $z \times T$: a reta de $T$
cai a $9{,}8$\,K/km, a de $T_d$ cai a $\sim 1{,}8$\,K/km; onde se cruzam,
nasce a nuvem. Marcar ``$z_{LCL}$'' no cruzamento.}

\begin{formadif}[por que as retas se cruzam]
Na subida seca, $T$ cai com $\Gamma_d = 9{,}8$\,K/km
\javimos{3}{gradiente adiabático seco}. O ponto de orvalho também cai
($P$ diminui $\Rightarrow$ $e$ diminui com $r$ constante), mas devagar:
$dT_d/dz \simeq -1{,}8$\,K/km (dedução no T4). As curvas se aproximam a
$8{,}0$\,K/km:
\[ z_{LCL} = \frac{T - T_d}{8{,}0\ \mathrm{K/km}}
 = 0{,}125\ \mathrm{km/K}\times (T - T_d). \]
\end{formadif}

\begin{formalivro}[Stull eq.~4.16]
\[ z_{LCL} = a\,(T - T_d),\qquad a = 0{,}125\ \mathrm{km\,°C^{-1}}. \]
\end{formalivro}

\begin{noquadro}[medir umidade olhando pela janela]
base plana dos cumulus $=$ LCL\\[2pt]
$T = 30\,°$C, $T_d = 20\,°$C $\Rightarrow$ base a
$0{,}125\times10 = 1{,}25$\,km\\[2pt]
quanto mais seco o ar, mais alta a base.
\end{noquadro}

No diagrama termodinâmico, o LCL é a regra de Normand: suba a adiabática
seca por $T$ e a iso-umidade por $T_d$; o cruzamento é o LCL
\veremos{5}{regra de Normand no Skew-$T$}. No notebook (Caso~3)
comparamos a regra linear com o cálculo exato do MetPy (diferença
$<5\%$).

\section{Temperatura de bulbo úmido}

Evaporação forçada esfria o termômetro até $T_w$
($T_d \le T_w \le T$, com igualdades só na saturação):

\begin{formalivro}[equação psicrométrica]
\[ r = r_s(T_w) - \gamma\,(T - T_w),\qquad
\gamma = \frac{C_p}{L_v} \simeq 0{,}40\ \mathrm{g\,kg^{-1}\,K^{-1}}
\ \ (\text{a } P = 100\ \mathrm{kPa}). \]
Medir $T$ e $T_w$ com dois termômetros baratos $\Rightarrow$ obter $r$.
É o balanço do Cap.~3 em miniatura: o calor sensível cedido pelo ar paga o
calor latente da evaporação \javimos{3}{razão de Bowen}.
\end{formalivro}

$T_w$ é também o limite fisiológico: bulbo úmido acima de $\sim 35\,°$C
impede o resfriamento por suor — o limite de habitabilidade em ondas de
calor úmidas \veremos{21}{extremos de calor}. E é $T_w$ (não $T$!) que
decide chuva vs.\ neve, porque o floco que cai evapora e se resfria
\veremos{7}{fase da precipitação}.

\section{Água total e água precipitável}\label{sec:total}

Água total da parcela: $r_T = r + r_L + r_i$ (vapor + líquido + gelo),
conservada em nuvem fechada sem chuva — a versão úmida da etiqueta
$(\theta, r)$ \javimos{4}{razão de mistura, \S\ref{sec:variaveis}}.

\begin{formadif}[quanta chuva cabe no céu]
Somando o vapor de toda a coluna (a mesma integral de massa do Cap.~1
\javimos{1}{massa da coluna}):
\[ W = \frac{1}{\rho_{liq}\,g}\int_{0}^{P_{sfc}} r\,dP
\qquad \text{(altura de água líquida, em mm).} \]
\end{formadif}

\begin{noquadro}[valores de $W$ para calibrar]
polar seco: $\sim5$\,mm \quad subtropical: 20--30\,mm \quad tropical
úmido: 50--70\,mm\\[2pt]
uma tempestade que ``chove 80\,mm'' importou vapor por
\emph{convergência} — não esvaziou só a própria coluna.
\end{noquadro}

``Rios atmosféricos'' — filetes de 400\,km de largura com $W$ alto e vento
forte — transportam mais água que o Amazonas
\veremos{13}{rios atmosféricos}. Cálculo de $W$ com MetPy no Caso~4.

\section{Balanço euleriano de água}\label{sec:balanco}

\quadro{Mesmo cubo euleriano do Cap.~3 no quadro, agora contabilizando
$r_T$: advecção, fluxo turbulento, fontes e sumidouros.}

\begin{formalivro}[balanço de umidade, forma simplificada]
\[ \frac{\Delta r_T}{\Delta t} \simeq
 -\,U\frac{\Delta r_T}{\Delta x}
 \;-\;\frac{\Delta F_{z,turb}}{\Delta z}
 \;+\;\text{(fontes: evaporação; sumidouros: precipitação)}. \]
Evaporação de superfície (fórmula \emph{bulk}, par do $F_H$
\javimos{3}{fórmula bulk}):
$F_{water} = C_E\,M\,(q_{sfc} - q_{ar})$, $C_E \sim 2\times10^{-3}$.
\end{formalivro}

Fechamos um ciclo de três capítulos: radiação fornece $F^*$
\javimos{2}{balanço radiativo}; Bowen decide quanto vira $F_E$
\javimos{3}{balanço de superfície}; o vapor sobe, condensa no LCL (este
capítulo) — e o Cap.~5 decide se a nuvem cresce
\veremos{5}{estabilidade}.

\section{Instrumentos}

Psicrômetro (par seco/úmido), higrômetro capacitivo (estações automáticas
e celulares), espelho resfriado ($T_d$ direto, padrão de calibração),
GPS-meteorologia (o atraso úmido do sinal $\propto W$ — cada antena GNSS é
um pluviômetro de vapor) \veremos{8}{sondagem remota}.

\section{Síntese da aula}

\begin{noquadro}[mapa final --- canto do quadro]
\[ \frac{de_s}{dT} = \frac{L_v e_s}{\Re_v T^2}
 \;\to\; e_s(T)\ \text{(dobra a cada 10\,K)}
 \;\to\; r = \frac{\varepsilon e}{P-e} \]
\[ z_{LCL} = 0{,}125\,(T-T_d) \qquad
   W = \frac{1}{\rho_{liq}\, g}\!\int\! r\,dP \]
\end{noquadro}

\section{Exercícios complementares}\label{sec:exercicios}

Soluções (professor) em \texttt{cap04\_solucoes.ipynb}.

\subsection*{Teóricos}

\begin{enumerate}[label=\textbf{T\arabic*.},itemsep=4pt]
  \item A partir de $de_s/dT = L_v e_s/(\Re_v T^2)$, integre com $L_v$
        constante e obtenha a fórmula do livro. Estime o erro de assumir
        $L_v$ constante sabendo que $L_v(T) \simeq 2{,}501 - 0{,}00237\,
        (T-273{,}15)$\,MJ/kg.
  \item Mostre que $e_s$ cresce $\sim$7\,\%/K perto de 300\,K
        (calcule $d\ln e_s/dT$) e conecte com a intensificação de chuvas
        extremas num clima 2\,K mais quente.
  \item Deduza $r = \varepsilon e/(P-e)$ a partir das leis de gás ideal
        para vapor e ar seco.
  \item Justifique $z_{LCL} = 0{,}125\,(T-T_d)$: mostre que
        $dT_d/dz \simeq -\Re_v T_d^2\,g/(L_v \Re_d T)$ numa subida com $r$
        constante e avalie numericamente.
  \item Derive $W = (1/\rho_{liq}g)\int r\,dP$ a partir da massa de vapor
        numa coluna (use o resultado T4 do Cap.~1) e calcule $W$ para uma
        atmosfera idealizada com $r(P) = r_0 (P/P_0)^3$, $r_0 = 20$\,g/kg.
\end{enumerate}

\subsection*{Numéricos (no notebook)}

\begin{enumerate}[label=\textbf{N\arabic*.},itemsep=4pt]
  \item Integre a EDO de Clausius--Clapeyron com $L_v(T)$ variável e
        compare com a fórmula de $L_v$ constante entre $-40$ e $+40\,°$C;
        trace o erro relativo.
  \item Psicrômetro digital: dado $(T, T_w, P)$, resolva a equação
        psicrométrica para $r$ e UR (com \texttt{scipy.optimize.brentq}) e
        valide contra o MetPy.
  \item Para a sondagem tropical do curso, calcule o LCL exato (MetPy) e
        pela regra dos 125\,m/°C; estime a base de nuvem ao longo de um dia
        em que $T$ sobe de 24 a 33\,°C com $T_d = 21\,°$C fixo.
  \item Água precipitável: compare $W$ da sondagem tropical com uma versão
        resfriada em 10\,K com a mesma UR; verifique a taxa de
        $\sim$7\,\%/K.
\end{enumerate}

\vspace{1em}
\noindent\rule{\linewidth}{0.4pt}\\
{\scriptsize Material derivado de R.~Stull, \emph{Practical Meteorology}
(CC BY-NC-SA 4.0); este documento herda a mesma licença.}

\end{document}
